\section{Variational representation and affine position}\label{sec:position}

We next represent a support functional by a vector field of the form
\eqref{eq:vector-field}, and choose coordinates in which the associated
norm power has no degree-two component. Throughout this section, $k$ has the value in
\eqref{eq:k-choice}, and $\eta$ is a finite positive Borel measure with
bounded support spanning $\R^n$. Neither symmetry of $\eta$ nor the
functional condition \eqref{eq:functional-condition} is needed here.

Let $A$ be a nonzero symmetric positive semidefinite matrix, and write
$V=\range(A)$. For an origin-symmetric convex body $L$ in $V$, with
interior relative to $V$, put
\[
 \tau_A(L)=\int h_L(Ax)\,d\eta(x).
\]
We extend its gauge to $\R^n$ by $g=g_L\circ P_V$, where $P_V$ denotes
orthogonal projection onto $V$. This extension is a seminorm, with
kernel $V^\perp$. Define
\begin{equation}\label{eq:variational-objective}
 J(g)=
 \begin{cases}
  \E\log g,&k=1,\\
  \E g,&k=2.
 \end{cases}
\end{equation}
The logarithm is integrated outside the kernel, which is a spherical
null set.
When $A$ is positive definite, these objectives are chosen so that their
Wulff first variations involve $g^{k-1}\nabla g=X_{g^k}$, the vector
field in \eqref{eq:norm-vector-field}.

This variational construction belongs to the framework for dual
curvature measures. The optimizer-to-measure argument corresponds to
the $L_1$ dual Minkowski problem with dual parameter $q=1-k$:
for $k=2$, compare \cite[Section~3.1, Lemmas~3.1--3.3]{HuangZhao2018};
for $k=1$ and full-dimensional bodies, minimization under $\tau_A(L)=1$
is equivalent, by polarity, to the $p=1$ entropy maximization in
\cite[Section~5, Lemmas~5.1--5.3]{HuangLutwakYangZhang2018}.
Wang and Zhou established uniqueness and continuity for the relevant
full-dimensional dual Minkowski problems
\cite[arXiv version, Theorems~1.1--1.3]{WangZhou2026}.
Here we prove the representation in our normalization, together with
uniqueness and continuity through changes of rank, where the gauges
become seminorms. This continuity is required to select affine
coordinates. The measures may have atoms, and all
variations remain valid for nonsmooth minimizing bodies.

\begin{lemma}\label{lem:minimizer}
For every nonzero symmetric positive semidefinite $A$, there is a unique
origin-symmetric convex body $L_A$ in $V=\range(A)$ that minimizes
\eqref{eq:variational-objective} subject to $\tau_A(L_A)=1$.
Denote its extended gauge by $g_A$.
\end{lemma}

\begin{proof}
The spanning hypothesis gives a norm on $V$:
\[
 q_A(y)=\int |y\cdot Ax|\,d\eta(x).
\]
Indeed, $q_A(y)=0$ implies that $Ay$ is orthogonal to $\supp\eta$,
and hence $Ay=0$; because $y\in V$, this forces $y=0$.
For any normalized competitor $L$ and $y\in L$, symmetry gives
$q_A(y)\le\tau_A(L)=1$. Thus all normalized competitors have a common
outer radius in $V$.

We also need a bound on their gauges along a minimizing sequence.
For any nonzero seminorm $g$, let $R=\max_{\Sph}g$. By separation,
$g$ is the support function of the compact symmetric set
$\{z:z\cdot y\le g(y)\text{ for every }y\in\R^n\}$.
This set has maximal radius $R$: if $Re$ is a point of maximal norm
in it, then
\begin{equation}\label{eq:seminorm-coordinate-bound}
 g(u)\ge R|u\cdot e|\qquad(u\in\Sph).
\end{equation}
Consequently,
\[
 J(g)\ge
 \begin{cases}
  \log R+\E\log|u_1|,&k=1,\\
  Rb_n,&k=2.
 \end{cases}
\]
Here $\E|\log|u_1||^2<\infty$: the first coordinate has a density
proportional to $(1-t^2)^{(n-3)/2}$ on $(-1,1)$, and
$\int_0^{1/2}|\log t|^2\,dt<\infty$. In particular every nonzero
seminorm has a finite logarithmic objective.

Set $b_A=\int|Ax|\,d\eta(x)>0$. The ball of radius $1/b_A$ in $V$
is a normalized competitor. Comparison with this ball bounds $J$
above along a minimizing sequence, so the preceding inequalities bound
$R$ above. The bodies therefore contain a common inner ball in $V$,
as well as lying in a common outer ball. Their gauges are uniformly
bounded and equicontinuous on the unit sphere of $V$. A subsequence
converges uniformly to the gauge of a body with the same two ball
bounds. The support functions converge uniformly as well, so the
normalization passes to the limit. The objective also passes to the
limit: in the logarithmic case the extended gauges are uniformly
comparable to $|P_Vu|$, whose logarithm is integrable by
\eqref{eq:seminorm-coordinate-bound}. This proves existence.

For uniqueness, let $L_0,L_1$ be two minimizers and
$L_t=(1-t)L_0+tL_1$, where $0<t<1$. Its normalization remains one.
For $P_Vu\ne0$, write $\rho_i(u)=1/g_{L_i}(P_Vu)$.
Taking points on the ray through $P_Vu$ gives
\[
 \frac{1}{g_{L_t}(P_Vu)}
 \ge (1-t)\rho_0(u)+t\rho_1(u).
\]
Both $-\log\rho$ and $1/\rho$ are strictly convex decreasing
functions of $\rho>0$. Thus the objective of $L_t$ is at most the
interpolated objectives, and the inequality is strict wherever
$\rho_0\ne\rho_1$. Distinct bodies have different radial functions
on an open set of directions in $V$, which gives a set of positive
spherical measure after projection. They would therefore give a
strictly smaller objective, a contradiction.
\end{proof}

\begin{proposition}[Variational representation]\label{prop:representation}
Suppose $A$ is positive definite, and let $L$ be any positive dilate
of $L_A$, with gauge $g=g_L$. For every continuous even function
$\phi:\R^n\to\R$ that is positively homogeneous of degree one,
\begin{equation}\label{eq:representation}
 \int\phi(Ax)\,d\eta(x)
 =s\,\E\phi\bigl(X_{g^k}(u)\bigr),
 \qquad
 s=\frac{\tau_A(L)}{\E g^{k-1}}.
\end{equation}
Here, almost everywhere on the sphere,
\begin{equation}\label{eq:norm-power-field}
 X_{g^k}=g^ku+g^{k-1}\gradS g=g^{k-1}\nabla g.
\end{equation}
\end{proposition}

\begin{proof}
Normalization turns the minimization problem into minimizing
$J(\tau_A(L)g_L)$ over all symmetric bodies. Put
$\tau=\tau_A(L)$ and $I=\int\phi(Ax)\,d\eta(x)$. For small
positive or negative $t$, define the body
\[
 W_t=\bigl\{y\in\R^n:
       y\cdot v\le h_L(v)+t\phi(v)\text{ for all }v\in\Sph\bigr\},
 \qquad g_t=g_{W_t}.
\]
The constraint is positive for sufficiently small $|t|$, so $W_t$ is
a symmetric convex body. Although that constraint need not be a
support function, it bounds the actual support function above.
Therefore
\[
 \tau_A(W_t)\le\tau+tI.
\]
Both quantities are positive for small $|t|$. Since $J$ increases
under pointwise increase of a positive gauge, minimality gives
\begin{equation}\label{eq:wulff-minimum}
 J\bigl((\tau+tI)g_t\bigr)
 \ge J\bigl(\tau_A(W_t)g_t\bigr)
 \ge J(\tau g),
\end{equation}
with equality at $t=0$.

Lemma~\ref{lem:gauge-variation} gives
\[
 \left.\frac{d}{dt}\right|_{t=0}g_t(u)
 =-g(u)\phi(\nabla g(u))
\]
almost everywhere, with uniform bounds that permit differentiation
of the objective. The left side of \eqref{eq:wulff-minimum} has a
two-sided minimum at zero. Its derivative is consequently zero:
\[
 \frac{I}{\tau}=\E\phi(\nabla g)\quad(k=1),
 \qquad
 I\E g=\tau\E\bigl[g\phi(\nabla g)\bigr]\quad(k=2).
\]
Equivalently,
\[
 \frac{I}{\tau}
 =\frac{\E\bigl[g^{k-1}\phi(\nabla g)\bigr]}{\E g^{k-1}}.
\]
Euler's identity for the homogeneous gauge gives
\eqref{eq:norm-power-field}, and positive homogeneity of $\phi$
then gives \eqref{eq:representation}. The argument applies to both
signs of $t$ without imposing convexity on the varied constraint.
\end{proof}

To select a position we allow the matrix to approach the boundary of
the compact convex set
\begin{equation}\label{eq:matrix-domain}
 \mathcal A=\{A=A^t\ge0:\tr A=1\}.
\end{equation}
The minimizers extend continuously to this boundary, even though
their ambient unit sets may become cylinders.

\begin{lemma}\label{lem:matrix-continuity}
The map $A\mapsto g_A|_{\Sph}$ is continuous on $\mathcal A$ in
the uniform norm.
\end{lemma}

\begin{proof}
Let $A_j\to A$ in $\mathcal A$, and write
$V_j=\range(A_j)$, $g_j=g_{A_j}$, and
$b_j=\int|A_jx|\,d\eta(x)$. Then
$b_j\to b=\int|Ax|\,d\eta(x)>0$. The normalized ball in $V_j$
has extended gauge $b_j|P_{V_j}u|$, and hence objective at most
$\log b_j$ if $k=1$, and at most $b_j$ if $k=2$.
By \eqref{eq:seminorm-coordinate-bound}, this bounds
$R_j=\max_{\Sph}g_j$ uniformly above. The minimizing body contains
the ball of radius $1/R_j$ in $V_j$, so its normalization gives
$1\ge b_j/R_j$. Thus $R_j\ge b_j$, bounded away from zero.

The seminorms $g_j$ have uniformly bounded Lipschitz constants.
Every subsequence therefore has a further subsequence converging
uniformly on $\Sph$ to a nonzero seminorm $g$. Along such a
subsequence,
\begin{equation}\label{eq:objective-limit}
 J(g_j)\longrightarrow J(g).
\end{equation}
For $k=2$ this is immediate. For $k=1$, choose unit vectors $e_j$
with $g_j(u)\ge R_j|u\cdot e_j|$. The positive parts of $\log g_j$
are uniformly bounded, and their negative parts have uniformly
bounded squared integrals, by rotation invariance and
$\E|\log|u_1||^2<\infty$. These logarithms are therefore uniformly
integrable. They converge almost everywhere to $\log g$, since
the kernel of a nonzero seminorm is a proper subspace. Truncating
the logarithms below at $-M$ and then sending $M$ to infinity proves
\eqref{eq:objective-limit}; the discarded integrals are uniformly
$O(M^{-1})$ by the squared-integral bound.

We identify this subsequential limit. Set
$Z=\{g\le1\}$ and $Z_j=\{g_j\le1\}$, allowing infinite values
for their support functions. If $g(y)<1$, then $y\in Z_j$ for all
sufficiently large $j$. Testing the support functions against $y$
and then taking the supremum gives, for every $x$,
\begin{equation}\label{eq:cylinder-support-limit}
 h_Z(Ax)\le\liminf_j h_{Z_j}(A_jx).
\end{equation}
The supremum over $g(y)<1$ equals $h_Z$, by scaling towards zero.
Since $Z_j=L_{A_j}+V_j^\perp$ and $A_jx\in V_j$, its support
function at $A_jx$ equals that of $L_{A_j}$. All these support
functions are nonnegative, so Fatou's lemma yields
\begin{equation}\label{eq:cylinder-normalization}
 \int h_Z(Ax)\,d\eta(x)\le1.
\end{equation}

Let $V=\range(A)$ and $L_0=\overline{P_VZ}$. This projection is
bounded despite the possible unboundedness of $Z$. Indeed, for
every $z\in Z$, symmetry and \eqref{eq:cylinder-normalization} give
\[
 q_A(P_Vz)=\int|z\cdot Ax|\,d\eta(x)
 \le\int h_Z(Ax)\,d\eta(x)\le1.
\]
The norm $q_A$ bounds the radius in $V$. Moreover $Z$ contains a
Euclidean ball, so $L_0$ has interior in $V$. Its extended gauge
$\bar g=g_{L_0}\circ P_V$ satisfies $\bar g\le g$, and
\begin{equation}\label{eq:projected-normalization}
 0<c_0:=\tau_A(L_0)=\int h_Z(Ax)\,d\eta(x)\le1.
\end{equation}

For the opposite comparison, let $N$ be any full-dimensional
symmetric convex body in $\R^n$. The extended gauge of its
projection to $V_j$ is at most $g_N$, and this projected body's
normalization is $\tau_{A_j}(N)$. Minimality and monotonicity give
\[
 J(g_j)\le J\bigl(\tau_{A_j}(N)g_N\bigr).
\]
Passing to the limit, using \eqref{eq:objective-limit} and
continuity of $\tau_{A_j}(N)$, gives
\begin{equation}\label{eq:full-body-comparison}
 J(g)\le J\bigl(\tau_A(N)g_N\bigr).
\end{equation}
Take $N=N_R=L_A+RB_{V^\perp}$, the orthogonal product of the
minimizing body in $V$ and a ball of radius $R$ in $V^\perp$.
Then $\tau_A(N_R)=1$ and
\[
 g_{N_R}(u)=\max\{g_A(u),|P_{V^\perp}u|/R\}\downarrow g_A(u).
\]
For $k=1$, the integrable function $\log g_A$ is a lower bound
for these logarithms, and their positive parts are uniformly
bounded for $R\ge1$; hence the objectives converge. The case
$k=2$ is immediate. Thus \eqref{eq:full-body-comparison} implies
$J(g)\le J(g_A)$.

Finally, $c_0\bar g$ is the extended gauge of the normalized
body $L_0/c_0$. Since $c_0\bar g\le g$, we have
\[
 J(g_A)\le J(c_0\bar g)\le J(g)\le J(g_A).
\]
Strict increase of the objective integrands forces
$c_0\bar g=g$ almost everywhere, and then everywhere by continuity.
Uniqueness in Lemma~\ref{lem:minimizer} gives
$c_0\bar g=g_A$. Thus every subsequential limit is $g_A$, which
proves the asserted continuity.
\end{proof}

We use Brouwer's fixed-point theorem to complete the choice of
position. We record the compact-convex formulation and its short
reduction to the simplex theorem \cite[Satz~4, p.~115]{Brouwer1911};
the nearest-point inequality in the reduction will also be used below.

\begin{lemma}[A finite-dimensional fixed-point principle]\label{lem:fixed-point}
Every continuous self-map of a nonempty compact convex subset of a
finite-dimensional Euclidean space has a fixed point.
\end{lemma}

\begin{proof}
Let the convex set be $D$, and work in its affine hull, of dimension
$d$. The case $d=0$ is immediate. The nearest-point projection
$\pi$ onto $D$ exists by compactness and is unique by strict
convexity of squared distance along a segment between distinct
minimizers. Minimizing along segments gives
\[
 \langle x-\pi(x),z-\pi(x)\rangle\le0\qquad(z\in D).
\]
Applying this twice yields
\[
 |\pi(x)-\pi(y)|^2
 \le\langle x-y,\pi(x)-\pi(y)\rangle,
\]
so $\pi$ is continuous. If $T:D\to D$ is continuous, choose a
$d$-simplex $\Delta$ containing $D$ and set
$G=T\circ\pi:\Delta\to\Delta$.

Brouwer's theorem for a simplex gives $v=G(v)$.
Since $G(\Delta)\subset D$, this point lies in $D$, hence
$\pi(v)=v$ and $T(v)=v$.
\end{proof}

\begin{proposition}[Affine position]\label{prop:position}
There is a symmetric positive definite matrix $A$ with $\det A=1$
such that $g_A^k$ has zero component in $\mathcal H_2$.
\end{proposition}

\begin{proof}
On $\mathcal A$ define the symmetric traceless matrix
\begin{equation}\label{eq:position-field}
 \mathcal F(A)
 =\frac{\E[g_A(u)^k uu^t]}{\E g_A^k}-\frac1n\Id.
\end{equation}
It is continuous by Lemma~\ref{lem:matrix-continuity}; its
denominator is positive because $g_A$ is a nonzero seminorm.

The direction $-\mathcal F(A)$ points strictly into the positive
cone along every missing direction of a singular matrix. To compute this, let
$A$ be singular and $e$ a unit vector in $\Ker A$. For a standard
Gaussian $Y$ in $\R^n$, $g_A(Y)$ depends only on the projection
of $Y$ to $\range(A)$, and is independent of $Y\cdot e$.
Separating the Gaussian radial and angular variables therefore gives
\begin{align*}
 \frac{\E[g_A(u)^k(u\cdot e)^2]}{\E g_A^k}
 &=\frac{\mathbb E[g_A(Y)^k(Y\cdot e)^2]}
          {\mathbb E g_A(Y)^k}
   \frac{\mathbb E|Y|^k}{\mathbb E|Y|^{k+2}}\\
 &=\frac{1}{n+k}.
\end{align*}
The last radial moment ratio follows by integration by parts in
$\int_0^\infty r^{n+k-1}e^{-r^2/2}\,dr$.
It follows, by polarization on $\Ker A$, that
\begin{equation}\label{eq:kernel-field}
 \left.\mathcal F(A)\right|_{\Ker A}
 =-\frac{k}{n(n+k)}I_{\Ker A}.
\end{equation}
The range--kernel block of $\mathcal F(A)$ is zero: reflecting
all kernel coordinates preserves $g_A$ and reverses each mixed
product. Consequently $A-\epsilon\mathcal F(A)$ is positive
definite for sufficiently small $\epsilon>0$. On the kernel this
follows from \eqref{eq:kernel-field}; on the range it follows from
the positive definiteness of the restriction of $A$. Its trace is
one. The same perturbation is allowed at a positive definite $A$
for sufficiently small $\epsilon$.

Use the Euclidean inner product $\langle U,W\rangle=\tr(UW)$ on
symmetric matrices, and consider the continuous map
\[
 A\longmapsto
 \operatorname{proj}_{\mathcal A}\bigl(A-\mathcal F(A)\bigr).
\]
Lemma~\ref{lem:fixed-point} provides a fixed point $A$. The
nearest-point condition there is
\[
 \langle\mathcal F(A),B-A\rangle\ge0
 \qquad(B\in\mathcal A).
\]
We may take $B=A-\epsilon\mathcal F(A)$, as just shown. This
gives $-\epsilon\tr(\mathcal F(A)^2)\ge0$, hence
$\mathcal F(A)=0$. Formula~\eqref{eq:kernel-field} rules out
singularity, so $A$ is positive definite.

For every symmetric traceless matrix $H$, the vanishing field gives
\[
 \E\bigl[g_A(u)^k u^tHu\bigr]=0.
\]
These quadratics are exactly $\mathcal H_2$, since the Euclidean
Laplacian of $x^tHx$ is $2\tr H$. Thus the required degree-two
component vanishes.

Finally, for $a>0$ the correspondence $L\mapsto L/a$ between
normalized competitors for $A$ and $aA$ gives
\begin{equation}\label{eq:minimizer-scaling}
 g_{aA}=a g_A.
\end{equation}
Indeed, $\tau_{aA}(L/a)=\tau_A(L)$, while multiplying a gauge by
$a$ adds $\log a$ to the logarithmic objective and multiplies
the linear objective by $a$. The minimizers therefore correspond.
Taking $a=(\det A)^{-1/n}$ makes the determinant one and preserves
the vanishing degree-two component.
\end{proof}

Only one measure needs the position of Proposition~\ref{prop:position}.
For a second measure, Proposition~\ref{prop:representation} provides
the required vector field in the resulting coordinates without a
simultaneous position condition.
